\section{Osservazioni}

Un eventuale allungamento della corda quando viene posta in trazione può influenzare le misure sperimentali in diversi modi.  
Nel nostro caso si è osservato sperimentalmente che la corda, quando viene tesa sulla macchina, può allungarsi di circa:

\[
\Delta L \simeq 7.2\ \text{mm}
\]

Tale effetto comporta innanzitutto una variazione della lunghezza effettiva del tratto vibrante della corda. Poiché la frequenza delle onde stazionarie dipende inversamente dalla lunghezza:

\[
f_n = \frac{nv}{2L}
\]

un aumento di $L$ produce una diminuzione della frequenza di risonanza rispetto al valore teorico atteso assumendo la corda perfettamente rigida.

Inoltre, l’allungamento modifica anche la massa lineare della corda. Infatti, a massa totale costante, un aumento della lunghezza comporta una lieve diminuzione della densità lineare:

\[
\mu = \frac{m}{L}
\]

e quindi una variazione della velocità di propagazione:

\[
v = \sqrt{\frac{\tau}{\mu}}
\]

Questo effetto tende ad aumentare leggermente la velocità dell’onda e quindi la frequenza. Tuttavia, nel nostro caso, la variazione della lunghezza del tratto vibrante rappresenta l’effetto dominante.

L’allungamento della corda può quindi introdurre una sistematica differenza tra valori teorici e sperimentali delle frequenze misurate, specialmente per tensioni elevate, e costituisce una possibile sorgente di errore sistematico nell’esperienza.